\section{Верификация и оценка на решението}

След представянето на архитектурата, домейн модела, back-end логиката, интерфейса и operational средата е необходимо разработката да бъде оценена систематично. Целта на тази глава не е просто да повтори, че системата съдържа дадена функционалност, а да покаже доколко решението удовлетворява поставените изисквания, какви са неговите силни страни и къде се намират естествените му граници.

В настоящия случай верификацията е извършена предимно чрез аналитичен подход върху реалния код, маршрути, шаблони, конфигурация и взаимовръзки между модулите. Това означава, че оценката се базира на действително реализираната логика на приложението, а не само на намерения или проектни декларации. Тъй като проектът има учебно-приложен характер и липсва богата автоматизирана тестова база, именно този тип инженерно-аналитична верификация е най-подходящ за дипломната работа.

\subsection{Подход към верификацията}

За целите на настоящата оценка са използвани няколко взаимнодопълващи се подхода.

\subsubsection{Статичен анализ на сорс кода}

Първият подход е проследяване на действителната логика в контролерите, услугите, repository интерфейсите, entity моделите и шаблоните. Чрез него може да бъде установено дали отделните потребителски сценарии имат реална програмна опора, каква е последователността на операциите и къде са разположени ключовите проверки за валидност, достъп или коректност на данните.

\subsubsection{Сценарийна проверка на функционалността}

Вторият подход се базира на дефинираните по-рано потребителски сценарии. За всеки от тях се проверява дали в системата има наличен контролерен маршрут, service логика, подходящ шаблон и persistence поведение, които заедно реализират целевия резултат. По този начин се формира сценарийна матрица на верификация.

\subsubsection{Архитектурна и качествена оценка}

Третият подход е по-качествен и разглежда не само наличието на дадена функция, а начина, по който тя е реализирана. Тук се оценяват свойства като структурираност на кода, повторна използваемост, разширяемост, защита на чувствителни данни, consistency на интерфейса и operational дисциплина.

Съвкупността от тези три перспективи дава достатъчно убедителна рамка за оценка на проекта като завършена дипломна разработка.

\subsection{Матрица за функционална верификация}

В таблица \ref{tab:verification-matrix} е представена обобщена матрица за функционална верификация. Тя не замества автоматизираното тестване, но систематизира дали ключовите функции са реално представени в кода и какви са доказателствата за тяхното съществуване.

\begin{longtable}{L{1.9cm}L{3.2cm}L{4.0cm}L{4.7cm}}
\caption{Матрица за функционална верификация}\label{tab:verification-matrix}\\
\toprule
Сценарий & Проверявана функционалност & Основа за верификация & Извод \\
\midrule
\endfirsthead
\toprule
Сценарий & Проверявана функционалност & Основа за верификация & Извод \\
\midrule
\endhead
V-1 & Регистрация на потребител & Налични \texttt{/users/register} GET/POST маршрути, service логика за запис и mail изпращане & Функцията е реализирана последователно и включва persistence и e-mail стъпка \\
V-2 & Потвърждение на акаунт & Наличен маршрут \texttt{/users/\{userId\}/confirm-email} и service метод \texttt{confirm} & Потвърждението е реализирано в опростен, но работещ вид \\
V-3 & Вход и изход & Налични login маршрути, проверка на парола и invalidate на сесията & Функцията е налице, но е базирана на ръчно сесийно управление \\
V-4 & Редакция на профил & Налични profile маршрути, проверки на стара парола и matching на нови пароли & Функцията е реализирана с базова входна защита \\
V-5 & Добавяне на фирма & \texttt{CompanyController}, \texttt{AddCompanyBindingModel}, \texttt{CompanyServiceImpl} & Реализацията е ясна и логически свързана с потребителя \\
V-6 & Добавяне на продукт & \texttt{ProductController}, \texttt{ProductAddBindingModel}, \texttt{ProductServiceImpl} & Налични са всички необходими елементи за запис на продукт \\
V-7 & Каталог по категории & Категорийни маршрути, шаблон \texttt{products-catalog.html}, enum категории & Функцията е реализирана, макар и с филтриране на ниво service вместо repository \\
V-8 & Детайлен преглед и добавяне в количка & Маршрут \texttt{/products/\{id\}/details}, проверка на количество, update на количката & Реализацията е налична, но изолацията по потребител е непълна \\
V-9 & Checkout & \texttt{shopping-cart} POST маршрут и \texttt{HistoryServiceImpl.secureCheckout} & Функцията е реализирана чрез транзакционна координация на няколко операции \\
V-10 & История на поръчките & \texttt{/products/my-history} GET/POST маршрути и \texttt{HistoryRepository} & Налице е основен исторически модул с възможност за изчистване \\
V-11 & Локализация & \texttt{messages.properties}, \texttt{messages\_bg.properties}, \texttt{I18nConfig} & Поддръжката е реална, но частично непълна \\
V-12 & REST крайна точка & \texttt{MaxPriceRestController} и service метод за максимална цена & Налице е примерна API функционалност \\
\bottomrule
\end{longtable}

Матрицата показва, че всички основни функции, формулирани в изискванията, имат конкретни програмни реализации. Това е съществен резултат, защото доказва, че системата не е само архитектурно упражнение, а реален функционален софтуерен продукт. В същото време няколко от сценариите ясно подсказват и ограничителни фактори, които следва да бъдат разгледани по-нататък.

\subsection{Оценка по основни качествени характеристики}

Функционалната верификация е необходима, но не е достатъчна. Една система може да „има“ дадена функция, но да я реализира по начин, който затруднява поддръжката, сигурността или бъдещото развитие. Поради това е необходимо проектът да бъде оценен и през няколко класически качествени измерения.

\subsubsection{Архитектурна подреденост}

От гледна точка на архитектурната подреденост проектът се представя добре. Налице е ясно разграничение между контролери, service слой, repository интерфейси, entity модели и визуален слой. Този факт значително улеснява анализа на системата и същевременно показва, че разработката следва устойчив организационен модел. Използването на отделни binding, service и view модели е допълнително свидетелство за осъзнато проектиране, а не само за ad hoc имплементация.

Слабост в тази характеристика е, че част от правилата за достъп и контрол на сесията са дублирани в контролерите. Това не нарушава функционалността, но намалява степента на централизация на сигурността.

\subsubsection{Поддържаемост}

Поддържаемостта е на добро ниво за мащаба на проекта. Кодът е сравнително четим, класовете имат разпознаваеми роли, а пакетната структура е логична. Това означава, че нов разработчик би могъл сравнително бързо да се ориентира в системата и да проследи основните потоци на данните.

Наред с положителната оценка тук трябва да се отчетат и няколко недостатъка: дублиране в HTML шаблоните, не напълно последователно именуване на някои полета и наличие на локално твърдо кодирани стойности. Тези елементи не правят системата неподдържаема, но биха били подходящ обект на следващ рефакторинг.

\subsubsection{Използваемост}

Използваемостта на приложението е приемлива и дори добра за учебен проект. Основните функции са групирани логично, навигацията е ясно видима, а потребителят може да премине през целия жизнен цикъл на покупката без да губи контекст. Каталогът, страницата за детайли, количката и профилът са особено важни в това отношение.

Ограниченията тук са свързани предимно с front-end зрелостта: частична локализация, разнороден визуален стил между отделни екрани, повторение на навигационни блокове и смесване на различни версии на UI библиотеки. Това не пречи на базовата употреба, но ограничава професионалната polish на интерфейса.

\subsubsection{Сигурност}

По отношение на сигурността системата демонстрира няколко правилни решения и няколко сериозни компромиса. Сред положителните аспекти са хеширането на паролите и забраната за вход на непотвърден потребител. Това са важни стъпки, които често липсват в базови учебни проекти.

Съществуват обаче и съществени ограничения:

\begin{itemize}[leftmargin=*]
    \item достъпът се управлява чрез \texttt{HttpSession} и ръчни условни проверки, без централизирана security рамка;
    \item чувствителни конфигурационни данни присъстват в \texttt{application.properties};
    \item контролът на данните в количката не е строго изолиран по потребител;
    \item няма роля-базирана авторизация и разграничение между административни и обикновени операции.
\end{itemize}

Следователно може да се заключи, че сигурността на системата е достатъчна за учебно-приложна демонстрация, но не и за production внедряване без значително доработване.

\subsubsection{Разширяемост}

Архитектурата на проекта позволява сравнително лесно добавяне на нови екрани, нови маршрути и нови service операции. Добавянето на нова страница следва установения модел: нов контролерен маршрут, нов шаблон, при нужда нов service метод и repository взаимодействие. От тази гледна точка разширяемостта е добра.

Най-същественото ограничение е, че част от бизнес правилата са имплементирани по начин, който би затруднил директното преминаване към multi-user, role-aware и production-grade система. Например количката и поръчковата история биха изисквали по-детайлно домейн моделиране, а security логиката -- централизация.

\subsection{Оценка на реализираните технологични решения}

Изборът на Spring Boot, Spring MVC, Thymeleaf, Spring Data JPA и MySQL може да бъде оценен като много подходящ за целите на настоящата разработка. Този стек осигурява достатъчно инженерна дълбочина, за да се разглеждат важни архитектурни и operational въпроси, но същевременно не изисква прекомерна инфраструктурна сложност. Това е една от причините проектът да бъде методологично силен като дипломна работа.

ModelMapper и Hibernate Validator допринасят за четимостта и дисциплината на кода, а mail интеграцията повишава функционалната стойност на системата. От operational гледна точка обаче dependency конфигурацията би могла да бъде по-чиста, особено по отношение на дублираните или твърде свободно декларирани версии.

\subsection{Наблюдавани ограничения, установени при анализа}

В резултат от направения преглед могат да бъдат формулирани няколко конкретни ограничения, които имат значение за цялостната оценка на проекта:

\begin{enumerate}[leftmargin=*]
    \item \textbf{Опростен security модел} -- контролът на достъпа е достатъчен за демонстрация, но не и за сериозна експлоатация.
    \item \textbf{Непълна изолация на количката} -- част от операциите работят върху глобални колекции, а не върху строго потребителски обхват.
    \item \textbf{Опростен order/history модел} -- липсват отделни позиции на поръчка и по-богат статусен модел.
    \item \textbf{Частична локализация} -- интерфейсът е двуезичен само частично и не напълно консистентно.
    \item \textbf{Inline front-end организация} -- повторение на навигация, локални стилове и смесване на UI ресурси.
    \item \textbf{Hard-coded конфигурация} -- локални пароли и адреси в configuration файла.
    \item \textbf{Ограничено автоматизирано тестване} -- архитектурата предполага възможност за тестове, но те не са развити като отделен доказателствен слой.
\end{enumerate}

Тези ограничения са важни именно защото са установени върху реалния код. Те не подкопават академичната стойност на проекта; напротив, показват, че системата е анализирана критично и професионално.

\subsection{Обобщена оценка спрямо целта на дипломната работа}

Основната цел на дипломната работа беше да се проектира и разработи уеб базирана система за електронна търговия с компютърна и офис техника, която да реализира основните бизнес процеси около потребителски акаунт, фирми, продукти, каталог, количка, checkout и история. На базата на проведената верификация може да се заключи, че тази цел е постигната.

Приложението е цялостно, стартиращо и методологично подредено. То реализира реални потребителски сценарии и използва индустриално релевантен технологичен стек. Допълнително проектът е достатъчно добре структуриран, за да бъде убедително описан в академичен текст и да служи като основа за демонстрация, защита и бъдещо доразвитие.

В същото време системата остава в границите на учебно-приложна разработка. Тя не претендира да бъде production-ready търговска платформа, а по-скоро да покаже как с разумен набор от технологии и ясно архитектурно мислене може да бъде изграден функционален модел на онлайн магазин. Именно в това се състои и нейната най-силна инженерна стойност.

\subsection{Заключение към оценъчната глава}

Верификацията и оценката показват, че проектът съчетава две ключови качества: функционална завършеност и аналитична прозрачност. Налични са реални потоци за регистрация, продуктово управление, покупка и история. Налице е смислена архитектура, която прави системата разбираема и надграждаема. Едновременно с това съществуват ясно разпознаваеми ограничения, които очертават следващия естествен етап в развитието на проекта.

Това прави разработката подходяща за дипломна работа: тя е достатъчно завършена, за да демонстрира практически резултат, и достатъчно критично анализируема, за да обоснове технически и академични изводи. Следващата глава систематизира именно тези ограничения и формулира посока за бъдеща работа.
